\chapter{Zarządzanie projektem}

\section{Planowanie pracy zespołu i środki komunikacji}
\par Na etapie planowania przyjęto iteracyjny model realizacji prac oparty na podejściu Agile, uzupełniony wybranymi praktykami SCRUM. Prace podzielono na krótkie iteracje o stałej długości: 2-tygodniowe \gls{sprint}y (z kilkoma wyjątkami). Ustalono także cykliczne, cotygodniowe spotkania, w ramach których definiowano zakres zadań, monitorowano postęp, konsultowano nowe pomysły albo napotkane problemy oraz dokonywano przegląd rezultatów. Każde spotkanie wspierane było agendą spotkania oraz \gls{meeting minutes} w celu efektywnego przebiegu obecnych jak i przyszłych spotkań. 

\section{Narzędzia wspomagające}
\par Zespół projektowy komunikował się głównie przy użyciu komunikatora Discord na wydzielonym do tego serwerze. Do przechowywania oraz kontroli wersji kodu źródłowego systemu użyty został serwis GitHub z wydzielonym repozytroium dostępnym pod adresem url: \\ https://github.com/Labbyn/labbyn . W celu przechowywania wszelkich wyprodukowanych podczas prac dokumentów, użyty został współdzielony Dysk Google. Do rejestrowania czasu pracy oraz zarządzaniu przydzielonymi zadaniami użyto serwisu Jira.

\section{Harmonogram Prac}
\par Ze względu na ograniczony czas na wyprodukowanie systemu, prace zostały podzielone na 5 etapów z umownymi datami rozpoczęcia i ukończenia etapu:
\begin{enumerate}
    \item Koncept (1. Września 2025 - 1. Listopada 2025)
    \item Rozpoczęcie prac (1. Listopada 2025 - 30. Listopada 2025)
    \item \gls{MVP} (1. Grudnia 2025 - 28. Stycznia 2026)
    \item Wersja 1.0 (29. Stycznia 2026 - 1. Maja 2026)
    \item Testy i poprawki (1. Maja 2026 - 20. Czerwca 2026)
\end{enumerate}
\par W celu najbardziej efektywnego postępu prac, zespół zdecydował się na poniższy proces tworzenia zadań oraz przyrostów:
\begin{itemize}
    \item Przyrosty będą opierały się na stworzonych i zaplanowanych zadaniach.
    \item Każde zadanie musi mieć przypisany \gls{Epik}.
    \item \gls{Epik}i będą tworzone na podstawie wymagań gdzie: 1 wymaganie = \gls{Epik}.
    \item \gls{Epik}i zostaną przypisane do konkretnego etapu w celu zaplanowania przyszłych prac.
\end{itemize}